% !TeX program = xelatex
% SPDX-License-Identifier: MIT
% Copyright (c) 2026 Ziyu Peng and 刘毅涵
% Author: https://github.com/pzy54154631-hub
% Author: https://pzy54154631-hub.github.io/liuyihan/
% 所有系数、标准误、样本量、R^2均为虚构排版占位值。
% 本文件不对应真实数据或实际回归，不可作统计推断。
\documentclass[UTF8,a4paper,fontset=fandol]{ctexart}
\usepackage[margin=2.5cm]{geometry}
\usepackage{amsmath,amssymb,booktabs,array}
\usepackage[hidelinks]{hyperref}
\title{计量经济学：结果表的组织}
\author{\href{https://github.com/pzy54154631-hub}{\textit{Ziyu Peng}}, University of Colorado Boulder,
  \\[0.1em]College of Arts and Sciences, Statistics and Data Science
  \\[0.1em]Boulder, CO 80309, USA
  \\[0.1em]\href{mailto:Ziyu.PengSr@colorado.edu}{Ziyu.PengSr@colorado.edu}
  \\[0.35em]
  \href{https://pzy54154631-hub.github.io/liuyihan/}{刘毅涵}，暨南大学经济学院}
\date{}
\begin{document}
\maketitle
\section{先写清楚模型}
考虑用于说明表格结构的横截面回归：
\begin{align}
 c_i&=\beta_0+\beta_1y_i+u_i,\label{eq:m1}\\
 c_i&=\beta_0+\beta_1y_i+\beta_2s_i+u_i.\label{eq:m2}
\end{align}
其中 $c_i$ 表示月消费（千元），$y_i$ 表示月收入（千元），
$s_i$ 表示家庭人数。每行观察值对应一个家庭。
\textbf{以下数值全部是虚构的排版占位值；没有实际运行回归，
也没有对应的数据集，不用于估计、检验或现实解释。}

\begin{table}[htbp]
\centering
\caption{结果表版式练习：虚构占位值}
\label{tab:regression}
\begin{tabular}{lrr}
\toprule
 & \multicolumn{2}{c}{因变量：月消费（千元）}\\
\cmidrule(lr){2-3}
 & \multicolumn{1}{c}{模型（1）}
 & \multicolumn{1}{c}{模型（2）}\\
\midrule
月收入（千元） & 0.620 & 0.580\\
 & (0.080) & (0.090)\\
家庭人数 & \multicolumn{1}{c}{—} & 0.350\\
 & & (0.120)\\
常数项 & 0.800 & 0.300\\
 & (0.250) & (0.310)\\
\midrule
观察值数量 & 120 & 120\\
$R^2$ & 0.410 & 0.460\\
\bottomrule
\end{tabular}
\par\smallskip
\begin{minipage}{0.88\linewidth}
\footnotesize
注：系数下方括号用于展示标准误应放置的位置。
所有系数、标准误、观察值数量和 $R^2$ 均为自设占位值；
破折号表示未纳入该变量。示例不标显著性星号。
实际研究须报告数据来源、样本筛选、估计方法及标准误类型；
如使用聚类标准误，还须说明聚类层级。
\end{minipage}
\end{table}

表~\ref{tab:regression} 练习了表头、数值、模型信息及表注的分工。
排版程序不会替代统计软件验证表内数字。
\end{document}
